
\section{Similitud técnica y régimen jurídico}
Un derivado militar puede compartir componentes con un tipo civil. La semejanza ayuda a comprender producto, pero no decide autoridad, aprobación ni reconocimiento. El artículo 2 del 2018/1139 delimita ámbito y exclusiones para actividades estatales concretas; no se convierte en exclusión general de toda aeronave pública de cualquier norma civil.

El RD 866/2015 aprueba el Reglamento de Aeronavegabilidad de la Defensa. Su uso exige identificar producto, actividad y disposiciones pertinentes. BOE-A-2022-21487 es una convocatoria de empleo público y no proporciona esa base.
\section{Armonización y reconocimiento}
MAWA es foro de armonización en el contexto de EDA. EMAR no se convierte por mera publicación en reglamento directamente obligatorio e idéntico para todos los Estados. Se comprueba implementación nacional y reconocimiento.

El reconocimiento tiene alcance y condiciones. Una aprobación civil o militar extranjera no es pase universal para cualquier tarea. El expediente identifica documento de aceptación, productos, funciones y límites. Si no consta, queda pendiente sin confundirlo con incapacidad técnica del taller.
\Needspace{.48\textheight}
\section{Configuración y misión}
Perfiles de misión, modificaciones, equipos y ambiente pueden diferir del derivado civil. Contadores y cargas responden a datos aplicables, sin heredar automáticamente todos los límites ni eliminar los que permanecen. Acceso a información también requiere control dentro del régimen.

Disciplina técnica y trazabilidad mantienen importancia en ambos ámbitos. Cambia el marco que asigna responsabilidades y reconoce evidencia; no existe un procedimiento único deducible para todos los productos militares.
\repaso{¿Puede un manual civil aportar información a un estudio militar?}{Puede, con alcance y configuración identificados. Su uso como dato aplicable y la certificación requieren el marco y aprobaciones correspondientes. La referencia por sí sola no crea reconocimiento.}
